\section{The matrix zeta formulation}
This section identifies exactly what the determinant generating function says about the level condition.
For commuting indeterminates $x_1,\ldots,x_n$, put $X=\operatorname{diag}(x_1,\ldots,x_n)$ and
\[
 Z_A(x)=\det(I-XA)^{-1}\in\C[[x_1,\ldots,x_n]].
\]
The inverse exists because the determinant has constant term one.
The squarefree case of MacMahon's master theorem, in the form of
\citet[Corollary~2.2, p.~608]{GJMaster}, states that
\[
 [x_1\cdots x_n]Z_A(x)=\per(A).
\]
Here it has the following precise consequence.

\begin{proposition}\label{prop:zeta}
For a complex linear space $L\subseteq\operatorname{Mat}_n(\C)$, the level condition is equivalent to
\[
 [x_1\cdots x_n]\bigl(Z_{J+t}(x)-Z_J(x)\bigr)=0
 \quad\text{for every }t\in L.
\]
After a linear parametrization of $L$, the homogeneous parts of this equation are exactly
$(n-k)!S_k|_L$, for $1\leq k\leq n$.
\end{proposition}
\begin{proof}
Apply the quoted coefficient identity to $A=J+t$ and $A=J$ and then apply
Proposition~\ref{prop:minc}. A linear parametrization preserves total degree in the parameters, so its degree-$k$ part is precisely the indicated multiple of $S_k|_L$.
\end{proof}

The formulation is therefore an equality of polynomial conditions, including their homogeneous degrees. A bound on the dimension of their common linear subspaces requires a further property of this generating function. The coefficient identity itself gives no additional vanishing orders. For example, at order two its value is
\[
 a+b+c+d+ad+bc
 \quad\text{for }t=\begin{pmatrix}a&b\\c&d\end{pmatrix},
\]
so the required homogeneous equations remain $a+b+c+d=0$ and $ad+bc=0$.
